% Første linje er dokumentklassen. Brug *altid* memoir!
\documentclass[a4paper,oneside,article]{memoir}

% Herunder indlæses forskellige pakker
\usepackage[utf8]{inputenc}  % Korrekt håndtering af æ, ø og å
\usepackage[T1]{fontenc}  % Korrekt håndtering af æ, ø og å
\usepackage{microtype}  % Typografisk magi! Giver bl.a. pænere orddeling
\usepackage{graphicx}  % Gør det muligt at indsætte billeder
\usepackage{amsmath}  % Giver adgang til uundværlige matematikting
\usepackage{mathtools}  % Retter ting i ams og giver ekstra muligheder
\usepackage[danish]{babel}  % Danske betegnelser og orddeling
\usepackage{float} %Mere kontrol over billede-placering
\renewcommand{\danishhyphenmins}{22}  % Bedre dansk orddeling
\usepackage[exponent-product=\ensuremath{\cdot}]{siunitx} % Sientific notation for forskellige ting, med en regel om at den skal bruge cdot istedet for times. 
\usepackage{derivative} % Til at skrive differentialligninger 
\usepackage{xcolor} % For at få flere tekst farver 
\usepackage[colorlinks=true,pdftex]{hyperref} % For at få refrences man kan trykke på, slet [] hvis man vil have kasser i stedet for farvet tekst.
\usepackage{pdfpages} % Til at indsætte andre pdf'er
\usepackage{lmodern} % Så man kan compile i TeXworks
\usepackage{alphabeta} % Til at skrive græske bogstaver
\usepackage{subfig} % For at have to billeder ved siden af hinanden
\usepackage{amssymb} % For \mathbb


% Pænere toc indents 
\usepackage{tocloft}
\cftsetindents{section}{1em}{2em}
\cftsetindents{subsection}{4em}{0em}


% Indholdet i dit dokument starter her!
\begin{document}
\author{Mikkel Moth Billing \\ 202307989}
\title{Lineær Algebra}
\date{\today}
\maketitle   % Indsætter overskriften
\setcounter{tocdepth}{2}
\setcounter{secnumdepth}{1}
\hypersetup{linkcolor=black} % Laver tableofcontents sort.  
\tableofcontents
\hypersetup{linkcolor=red} % og så alle andre hyperrefs røde. 

%Custom Commands
\newcommand{\opg}[3]{
\hyperref[asc:#1]{#2}

\textcolor{gray}{#3} \plainbreak{1}
} %Til opgavebeskrivelser

\newcommand{\ti}[1]{\cdot 10^{#1}} % Let: gange 10^x
\newcommand{\e}[1]{\emph{e}^{#1}} % Pænere e^x
% Lette paranteser 
\newcommand{\br}[2][1]{%
    \ifcase#1
    \or \left( #2 \right) %1 eller ingenting
    \or \left[ #2 \right] %2
    \or \left\langle #2 \right\rangle %3
    \or \left\{ #2 \right\} %4
    \or \left| #2 \right| %5
    \fi
}
\newcommand{\bil}[1]{
\begin{figure}[H]
    \centering
    \includegraphics[width=0.85\linewidth]{Billeder/#1.png}
    \label{fig:#1}
\end{figure}
}

\newcommand{\pdf}[4]{
\includepdf[pages=#3,scale=.85,pagecommand={\subsection{#2 - \ref{ch:opg}}  \label{asc:#2}}]{PDF/#1.pdf}
\includepdf[pages=\the\numexpr 1 + #3\relax-#4, scale=.85]{PDF/#1.pdf}
}

\newcommand{\prop}[1]{
\hyperref[fig:#1]{(#1)}
}

% Lav ny pagestyle (kun odd fordi enkeltsidet)
\makepagestyle{Title}
\pagestyle{Title}
\makeoddhead{Title}{}{}{\footnotesize\color{gray}{\theauthor}}
\makeoddfoot{Title}{}{\footnotesize\color{gray}{\thepage{} / \thelastpage}}{}
% Også fod på siden med \maketitle
\makeoddfoot{plain}{}{\footnotesize\color{gray}{\thepage{} / \thelastpage}}{}
\newpage

%Skriv
\chapter{Generalt} \label{ch:gen}
\section{Matrixprodukt}
\bil{3.4}

\section{Løsningsmængde}
\bil{3.12}

\section{Injektiv}
\bil{6.10}
\section{Surjektiv} \label{sc:sur}
At en lineær transformation $L: V \rightarrow W$ er surjektiv betyder, at der for enhver vektor $w \in W$ findes en vektor $v\in V$ så $L(v)=w$. For at se om en lineær transformation er surjektiv, kan man bruge 
\bil{7.3}
Hvis rang(\emph{A})=dim(\emph{W}) er $L_A$ surjektiv. Hvis \emph{L} både er injektiv og surjektiv, er den en isomorfi. 

\hyperref[asc:Eksamen 2020 opg 1]{Eksempelopgave}

\section{Rang} \label{sc:rang}
\bil{7.37}
\bil{7.38}

\section{Similær}
\bil{12.21}

\chapter{Invertible matricer} \label{ch:inv}
\bil{4.1}
\textcolor{white}{singulær}
\bil{4.3}
\bil{4.4}
\bil{4.6}
\bil{7.4}
\bil{11.12}
En matrix med 2 ens søjler er ikke invertibel, og determinanten er derfor lig 0. 
\bil{8.25}
\hyperref[asc:Eksamen 2020 opg 3]{Eksempelopgave}

\chapter{Vektorrum} \label{ch:rum}
\section{Underrum} \label{sc:urum}
\bil{5.6}
\hyperref[asc:Reeksamen 2022 opg 1]{Eksempelopgave}
\bil{6.7}
\bil{7.16}
\hyperref[asc:Eksamen 2021 opg 1]{Eksempelopgave}

\section{Nulrum}
\emph{n} er antallet af søjler, \emph{r} er antallet af pivoter
\bil{7.30-33}
\bil{7.34}
\bil{7.36}
\hyperref[asc:Eksamen 2021 opg 1]{Eksempelopgave}

\section{Span}
\bil{5.12}
Hvis man skal vise at $V=\text{Span}(v_1,...,v_n)$.
Så skal man vise at for alle $A\in V$, så gælder $A \in \text{Span}(v_1,...,v_n)$, hvor $(v_1,...,v_n)$ er en basis. Dermed er $A\subseteq \text{Span}(v_1,...,v_n)$, og fra (5.12) har vi $A\supseteq \text{Span}(v_1,...,v_n)$.


\chapter{Lineære transformationer} \label{ch:trans}
For at vise at \emph{L} er en lineær transformation/operator, kan man bruge: 
\bil{6.1}
En anden måde, er hvis $L=L_A$ altså hvis man kan lave transformationen ved at gange med en matrix \emph{A}
\hyperref[asc:Eksamen 2023 opg 3]{Eksempelopgave}, \hyperref[asc:Reeksamen 2022 opg 3]{Eks.opg.}
\bil{6.20}
\bil{7.20}

\section{Invertible LT og isomorfi}
\bil{6.13}
\bil{6.14}
\bil{6.15}
\bil{6.16}
\bil{6.17}
For invertible lineære transformationer gælder: $(L\circ L^-1)(A)=A$, hvis dette er er opfyldt, er \emph{L} lineær isomorf. 

\hyperref[asc:Reeksamen 2023 opg 3]{Eksempelopgave}

\bil{8.25}

\chapter{Baser og lineær uafhængighed} \label{ch:baser}
\bil{7.4}
\bil{7.6}
\bil{7.9}
\bil{7.11}
\bil{7.12}
\bil{7.20}

\section{Baser for polynomier}
Hvis en opgave spørger efter en base for et n'te grads polynomium, kan eksempel 7.7(3) bruges: 
\bil{7.7(3)}
Man kan også sige at $\mathcal{V}=(1,X)$ er lineær uafhængig, da \emph{X} ikke er en skalar. 

\chapter{Matrixrepræsentationer} \label{ch:rep}
\bil{8.3}
\textcolor{white}{Koordinattransformationsmatrix}
\bil{8.6}
\bil{8.8}
\bil{8.10(2)}
\hyperref[asc:Reeksamen 2021 opg 1]{Eksempelopgave}
\bil{8.18}
Denne korollar kan også bruges, hvis det er to baser for det samme vektorrum, så er de mærkede den samme base, og de umærkede den anden. 
\begin{align*}
    \prescript{}{\mathcal{W}}[L]_\mathcal{W} &= \prescript{}{\mathcal{W}}[\Box]_\mathcal{V}\cdot \prescript{}{\mathcal{V}}[L]_\mathcal{V}\cdot \prescript{}{\mathcal{W}}[\Box]_\mathcal{V} \\
    \prescript{}{\mathcal{W}}[L]_\mathcal{W} &= \prescript{}{\mathcal{W}}[\Box]_\mathcal{V}\cdot \prescript{}{\mathcal{V}}[L]_\mathcal{W}
\end{align*}
\bil{8.21}
Eksempel 8.22.2: 
\bil{8.22.2}
\bil{8.23} 



\chapter{Indre produkt} \label{ch:indre} 
\bil{9.1}
For at vise at noget er et indre produkt, skal disse 4 krav vises. 

\hyperref[asc:Reeksamen 2023 opg 2]{Eksempelopgave}, \hyperref[asc:Eksamen 2022 opg 3]{Eks.opg.}

\bil{9.3.1}
\[\bar{w}^T\bar{v}=\br[3]{v,w}\]

\section{Norm}
\bil{9.6-7}
Fra TØ8: 
Et normeret K-vektorrum er et K-vektorrum V , hvorpå der er defineret en afbildning
\[||\cdot|| : V \rightarrow [0, \infty [,\]
opfyldende (for $\Bar{v}, \Bar{w} \in V $)

a) $||v|| = 0$ hvis og kun hvis $\bar{v} = 0$.

b) $||α \cdot \bar{v}|| = |α| \cdot ||\bar{v}||$ for $α \in \mathbb{K}$.

c) $||\bar{v} + \bar{w}|| \leq ||\bar{v}|| + ||\bar{w}||$.

\bil{9.17}

\section{Afstand}
\bil{9.8}
Afstandsbevarelse: 
\begin{align*}
    & L: V \rightarrow V \\
    & \forall \; \bar{v}, \bar{w} \in V \\
    ||L(\bar{v})- & L(\bar{w})|| = ||\bar{v}-\bar{w}||
\end{align*}


\chapter{Ortogonalitet} \label{ch:orto}
\bil{10.1}
Dette bruges til at teste om noget er en ortogonal/ortonormal mængde. 
For at regne fra ortogonal basis $(v_1,...,v_n)$ til ortonormal basis $(\mathcal{W})$: 
\[\mathcal{W} = \br{\frac{v_1}{||v_1||},...,\frac{v_n}{||v_n||}}\]

Bemærk: \[\br[3]{v_i,v_j}=0 \Leftrightarrow v_i \perp v_j\]
\bil{10.4}
Sammen med Prop. \prop{7.12}, kan man sige at $(v_1 ... v_n)$ er en ortonormal basis for $\mathbb{R}^n$
\section{Ortogonal projektion} \label{sc:op}
\bil{10.14}
Her er en måde at finde den ortogonale projektion \emph{p}. Her er $v_n$ en ortogonal mængde (f.eks den base fundet ved Gram-Schmidt) og \emph{v} er den ting, man vil finde projektionen af. 

\hyperref[asc:Eksamen 2023 opg 2]{Eksempelopgave}
\bil{10.18}

\section{Ortogonalt komplement} \label{sc:ok}
\bil{10.5}
Det ortogonale komplement er de vektorer der er ortogonale på alle vektorerene i \emph{W}. Der gælder også at: 
\bil{10.22}
\bil{10.41}
Hvis \emph{V} er 3-dimensionelt rum, og \emph{W} består udelukkende af vektorer i x-y planen, så består $W^\perp$ af alle vektorer i z-planen. 
\plainbreak{1}
For at finde en basis til det ortogonale komplement, kan vi fra \hyperref[fig:10.14]{(10.14)} se at \[\bar{h}=\bar{p}-\bar{v} \in W^\perp\] 

\hyperref[asc:Eksamen 2023 opg 2]{Eksempelopgave}

\section{Lineær isometri}
\bil{10.28}

\section{Gram-Schmidt processen} \label{sc:gs}
Husk at man skal bruge det indre produkt, der er opgivet i opgaven, når man finder normen. 
\bil{10.25}

\section{Unitær og ortogonal matricer}
\bil{10.31}
\bil{10.34}

\section{Mindste kvadrats løsning}
\bil{10.39}
\bil{10.40}
\bil{10.42}
\hyperref[asc:Eksamen 2020 opg 1]{Eksempelopgave}

\chapter{Determinanter} \label{ch:det}
\bil{11.9}
\bil{11.23}
\bil{11.12}
\bil{11.20}
\[det(A)\cdot det(A^{-1}) = det(AA^{-1})=det(I)=1\]

\section{Kofaktor og adjungeret}
\bil{11.16}
\bil{11.24-25}
\bil{11.27}
\hyperref[asc:Eksamen 2020 opg 4]{Eksempelopgave}

\chapter{Egenværdier} \label{ch:egen}
\bil{12.1}
$A\cdot \bar{v} = λ\cdot \bar{v}$ er vigtig til mange opgaver. Hvis $A-λI$ er singulær er $λ=0$ jf. TØ12 opg 4. 

\hyperref[asc:Eksamen 2023 opg 5]{Eksempelopgave}
\bil{12.5}

\section{Egenrum} \label{sc:erum}
\bil{12.10}
\bil{12.12}
Dimensionen af nulrummet kan let findes, da det er antallet af nulrækker i en matrix på RREF. 
\bil{12.13}

\section{Karakteristisk polynomium} \label{sc:poly}
Det karakteristiske polynomium for \emph{A} (eller $_{\mathcal{V}}[L]_{\mathcal{V}}$) findes som: 
\bil{12.15}
\bil{12.23}
\bil{12.24}
\bil{12.21}


\hyperref[asc:Eksamen 2023 opg 1]{Eksempelopgave}

\section{Algebraisk Multiplicitet} \label{sc:alg}
\bil{12.28}
$\text{Alg}_A(λ)$ findes ved, at man faktorisere det karakteristiske polynomium, og kigger på roden af de forskellige egenværdier. 
\bil{alg1}
En vigtig relation til den geometrike multiplicitet: 
\bil{12.30}
Dette kan bruges med \hyperref[fig:13.13]{(13.13)}, til at se om A er diagonaliserbar. 

\chapter{Diagonalisering} \label{ch:diag}
\bil{13.9}
\bil{13.10}
Gælder også for $L_A=A$. 
\bil{13.13}

\bil{13.5}
\hyperref[asc:Eksamen 2022 opg 1]{Eksempelopgave}
\bil{14.18} \ref{sc:adj}

\section{Ortonormal diagonalisering}
\bil{14.1}
\bil{14.4}
\bil{14.5}
\bil{14.22}

\section{Selvadjungeret} \label{sc:adj}
\bil{14.11}
\hyperref[asc:Eksamen 2022 opg 3]{Eksempelopgave}
\bil{14.15}
\bil{14.18} 

\hyperref[asc:Eksamen 2022 opg 3]{Eksempelopgave}

\chapter{Differentialligninger} \label{ch:diff}
\bil{15.2}
\section{Eksponentialfunktion} \label{sc:eks}
\bil{15.4}
For at vise at noget er en eksponentialfunktion, kan man vise at disse 3 krav er opfyldt. 

\hyperref[asc:Eksamen 2023 opg 5]{Eksempelopgave}
\bil{15.8}

\section{Lineære differentialligningssystemer} \label{sc:cock}
\bil{cock1}
\bil{15.16}
$\bar{v} = f(0)$, \hyperref[asc:Reeksamen 2023 opg 4]{Eksempelopgave}
\bil{15.18}
\bil{15.19}
\bil{15.20}
\hyperref[asc:Eksamen 2021 opg 4]{Eksempelopgave}

\newpage
\addtocontents{toc}{\protect\setcounter{tocdepth}{1}} 
\setcounter{tocdepth}{1}
\chapter{Opgaver} \label{ch:opg}

% 1
\opg{Eksamen 2023 opg 1}{Eksamen 2023 opg 1}{Karakteristisk polynomium, algebrarisk multiplicitet, diagonaliserbar}
\opg{Eksamen 2023 opg 2}{Eksamen 2023 opg 2}{Polynomier, ortogonal projektion, ortogonalt komplement}
\opg{Eksamen 2023 opg 3}{Eksamen 2023 opg 3}{Determinant, lineær transformation, basis, surjektiv}
\opg{Eksamen 2023 opg 4}{Eksamen 2023 opg 4}{Ortonormalbasis, matrixrepræsentation, selvadjungeret, ortonormalt diagonaliserbar}
\opg{Eksamen 2023 opg 5}{Eksamen 2023 opg 5}{Egenværdier, lineære differntialligningssystemer}
\plainbreak{1}

% 2
\opg{Reeksamen 2023 opg 1}{Reeksamen 2023 opg 1}{Matrixrepræsentationen, karakteristisk polynomium, egenværdi, diagonaliserbar}
\opg{Reeksamen 2023 opg 2}{Reeksamen 2023 opg 2}{Indre produkt, Gram-Schmidt, egenvektorer, ortonormalt diagonaliserbar, selvadjungeret}
\opg{Reeksamen 2023 opg 3}{Reeksamen 2023 opg 3}{Lineær isomorfi, underrum, ortogonal basis, ortogonalt komplement, projektion}
\opg{Reeksamen 2023 opg 4}{Reeksamen 2023 opg 4}{Eksponentielfunktion, lineære differentialligningssystem}
\plainbreak{1}

% 3
\opg{Eksamen 2022 opg 1}{Eksamen 2022 opg 1}{Egenværdier, diagonaliserbar, similær, determinant}
\opg{Eksamen 2022 opg 2}{Eksamen 2022 opg 2}{Ortogonal basis, egenvektor, ortonormalt diagonaliserbar, selvadjungeret, ortogonal projektion}
\opg{Eksamen 2022 opg 3}{Eksamen 2022 opg 3}{Ortonormal basis, selvadjungeret operator, indre produkt}
\opg{Eksamen 2022 opg 4}{Eksamen 2022 opg 4}{Differentialligninger, matrixrepræsentation, nulrum}
\plainbreak{1}

% 4
\opg{Reeksamen 2022 opg 1}{Reeksamen 2022 opg 1}{Vis underrum, lineær uafhængighed og dimension}
\opg{Reeksamen 2022 opg 2}{Reeksamen 2022 opg 2}{Basis, Gram-Schmidt, projektion}
\opg{Reeksamen 2022 opg 3}{Reeksamen 2022 opg 3}{polynomier, lineær operator, matrixrepræsentation, egenværdi, lineær isomorfi, diagonaliserbar}
\opg{Reeksamen 2022 opg 4}{Reeksamen 2022 opg 4}{Differentialligningssystem, lineært uafhængig, basis}
\plainbreak{1}

% 5 
\opg{Eksamen 2021 opg 1}{Eksamen 2021 opg 1}{Basis, søjlerum, nulrum, lineær ligningssystem}
\opg{Eksamen 2021 opg 2}{Eksamen 2021 opg 2}{Gram-Schmidt, ortonormal basis, multiplicitet, selvadjungeret}
\opg{Eksamen 2021 opg 3}{Eksamen 2021 opg 3}{Indre produkt, ortogonal projektion, udregninger}
\opg{Eksamen 2021 opg 4}{Eksamen 2021 opg 4}{Differentialligning basis, lineær kombination, lineær operator, matrixrepræsentation, isomorfi}
\plainbreak{1}

% 6
\opg{Reeksamen 2021 opg 1}{Reeksamen 2021 opg 1}{Koordinattransformationsmatrice, matrixrepræsentation, ortogonal basis}
\opg{Reeksamen 2021 opg 2}{Reeksamen 2021 opg 2}{Ortogonal mængde, rang, ortogonalt komplement/basis/projektion}
\opg{Reeksamen 2021 opg 3}{Reeksamen 2021 opg 3}{Basis, matrixrepræsentation, egenværdier, vis indre produkt, Gram-Schmidt}
\opg{Reeksamen 2021 opg 4}{Reeksamen 2021 opg 4}{Underrum, basis ved hjælp af Span og underrum, eksponentielfunktion, \hyperref[fig:15.4]{(15.4)}, differentialligningssystem}
\plainbreak{1}

% 7
\opg{Eksamen 2020 opg 1}{Eksamen 2020 opg 1}{Injektiv og ikke surjektiv, ingen løsninger til lineær ligningssystem, mindste kvadraters løsning, ortogonal projektion}
\opg{Eksamen 2020 opg 2}{Eksamen 2020 opg 2}{Integrale, selvadjungeret, ortonormalt diagonaliserbar, karakteristisk polynomium, egenværdi}
\opg{Eksamen 2020 opg 3}{Eksamen 2020 opg 3}{Matrixrepræsentation $_{\mathcal{E}}[L]_{\mathcal{V}}$, isomorfi, vis indre produkt, Gram-Schmidt med funktioner, ortogonal projektion}
\opg{Eksamen 2020 opg 4}{Eksamen 2020 opg 4}{Adjungeret, eksponentielfunktion, lineær differentialligningssystem med begyndelsesværdi}
\plainbreak{1}

% 8 
\opg{Reeksamen 2020 opg 1}{Reeksamen 2020 opg 1}{Indre produkt og kerne, injektiv, isomorfi, selvadjungeret, ortonormalt diagonaliserbar}
\opg{Reeksamen 2020 opg 2}{Reeksamen 2020 opg 2}{Søjlerum, mindste kvadraters løsning, nulrum, løsningsmængde, ortogonal projektion}
\opg{Reeksamen 2020 opg 3}{Reeksamen 2020 opg 3}{Indre produkt, lineært uafhængig, Gram-Schmidt, koordinattransformationsmatricen}
\opg{Reeksamen 2020 opg 4}{Reeksamen 2020 opg 4}{Egenvektorer, egenværdier, eksponentielfunktioner, differentialligningssystemer}
\plainbreak{1}

% 9
\opg{Eksamen 2019 opg 1}{Eksamen 2019 opg 1}{Basis med determinant, Gram-Schmidt}
\opg{Eksamen 2019 opg 2}{Eksamen 2019 opg 2}{Ikke invertibel, basis for søjlerum, koordinattransformationsmatrice}
\opg{Eksamen 2019 opg 3}{Eksamen 2019 opg 3}{Løsningsmængde, basis for nulrum, dimension af række/søjlerum, ortogonalt komplement}
\opg{Eksamen 2019 opg 4}{Eksamen 2019 opg 4}{Vis indre produkt, egenrum, egenværdi, selvadjungeret, diagonaliserbar, ortonormal basis, isometri}
\opg{Eksamen 2019 opg 5}{Eksamen 2019 opg 5}{Eksponentielfunktion,  differentialligningssystem}

\newpage
\subsection{TØ opgaver: }
\textcolor{white}{-}

Uge 2: \textcolor{gray}{ Matricer, RREF, udsagn og induktionbevis. }

Uge 3: \textcolor{gray}{ Matrixprodukt, invertible matricer og rækkeækvivalens. }

Uge 4: \textcolor{gray}{ Vektorrum, underrum, span og ligningssystemer uden løsning. }

Uge 5: \textcolor{gray}{ Billede, kerne, standardmatrixrepræsentation, nulrum, lineær transformation/operator, isomorfi, injektiv og surjektiv. }

Uge 6: \textcolor{gray}{ Span, lineær uafhængighed og basis. }

Uge 7: \textcolor{gray}{ Underrum, kvotientrum, basis for søjlerum, rækkerum og nulrum og koordinattransformationsmatrice. }

Uge 8: \textcolor{gray}{ Basis, koordinattransformationsmatrix, matrixrepræsentation, indre produkt og normeret vektorrum.}

Uge 9: \textcolor{gray}{ Indre produkt rum og underrum, projektion, ortogonal basis, ortogonal/ortonormal mængde, nulrum og spejlinger i hyperplaner. }

Uge 10: \textcolor{gray}{ Gram-Schmidt, mindste kvadrats løsning, ortogonal projektion, lineær/kvadratisk/kubisk approksimation, isomorfi, afstandsbevarende, polariseringsidentiteten og ortonormal basis}

Uge 11: \textcolor{gray}{Matrixrepræsentation, lineær approksimation, determinant, skævsymmetrisk og ortonormal basis.}

Uge 12: \textcolor{gray}{Ortogonal basis, egenværdier og isometri for cos og sin, $\text{det}\br{A^TA}$, egenværdi, egenværdi for projektionsafbildning.}

Uge 13: \textcolor{gray}{Induktionsbevis, karakteristisk polynomium, egenværdi, multipliciteter og egenværdier for polynomier.}

Uge 14: \textcolor{gray}{Egenværdi, ortonormalt diagonaliserbar, mindste kvadraters løsning, Gram-Schmidt, koordinattransformationsmatrix, selvadjungeret og komplekse funktioner.}

Uge 15: \textcolor{gray}{Reelle og komplekse differentialligningssystemer, eksponentialfunktion, }

%\begin{comment}
% 1
\pdf{23v}{Eksamen 2023 opg 1}{1}{2}
\pdf{23v}{Eksamen 2023 opg 2}{3}{4}
\pdf{23v}{Eksamen 2023 opg 3}{5}{6}
\pdf{23v}{Eksamen 2023 opg 4}{8}{9}
\pdf{23v}{Eksamen 2023 opg 5}{10}{12}

% 2
\pdf{23rv}{Reeksamen 2023 opg 1}{1}{3}
\pdf{23rv}{Reeksamen 2023 opg 2}{4}{6}
\pdf{23rv}{Reeksamen 2023 opg 3}{7}{10}
\pdf{23rv}{Reeksamen 2023 opg 4}{11}{13}


% 3 
\pdf{22v}{Eksamen 2022 opg 1}{1}{2}
\pdf{22v}{Eksamen 2022 opg 2}{3}{5}
\pdf{22v}{Eksamen 2022 opg 3}{6}{8}
\pdf{22v}{Eksamen 2022 opg 4}{9}{11}

% 4 
\pdf{22rv}{Reeksamen 2022 opg 1}{1}{4}
\pdf{22rv}{Reeksamen 2022 opg 2}{5}{7}
\pdf{22rv}{Reeksamen 2022 opg 3}{8}{10}
\pdf{22rv}{Reeksamen 2022 opg 4}{11}{13}

% 5
\pdf{21v}{Eksamen 2021 opg 1}{1}{3}
\pdf{21v}{Eksamen 2021 opg 2}{4}{7}
\pdf{21v}{Eksamen 2021 opg 3}{8}{11}
\pdf{21v}{Eksamen 2021 opg 4}{12}{14}

% 6
\pdf{21rv}{Reeksamen 2021 opg 1}{1}{3}
\pdf{21rv}{Reeksamen 2021 opg 2}{4}{6}
\pdf{21rv}{Reeksamen 2021 opg 3}{7}{12}
\pdf{21rv}{Reeksamen 2021 opg 4}{13}{16}

% 7 
\pdf{20v}{Eksamen 2020 opg 1}{1}{3}
\pdf{20v}{Eksamen 2020 opg 2}{4}{6}
\pdf{20v}{Eksamen 2020 opg 3}{7}{10}
\pdf{20v}{Eksamen 2020 opg 4}{11}{14}

% 8 
\pdf{20rv}{Reeksamen 2020 opg 1}{1}{3}
\pdf{20rv}{Reeksamen 2020 opg 2}{4}{7}
\pdf{20rv}{Reeksamen 2020 opg 3}{8}{11}
\pdf{20rv}{Reeksamen 2020 opg 4}{12}{14}

% 9 
\pdf{19v}{Eksamen 2019 opg 1}{1}{3}
\pdf{19v}{Eksamen 2019 opg 2}{4}{5}
\pdf{19v}{Eksamen 2019 opg 3}{6}{8}
\pdf{19v}{Eksamen 2019 opg 4}{9}{11}
\pdf{19v}{Eksamen 2019 opg 5}{12}{13}


\chapter{Appendix}
\pdf{Matnoter}{Inductionsbeviser}{53}{56}

%\end{comment}

\end{document} % ----------Dokumentet er nu slut! ------------